\Qhead{4.5}{Молекулярный генератор. Квантовые стандарты частоты (времени)}{Молекулярный генератор. Квантовые стандарты частоты (времени).}
\Subfile{4.5.1}{sections/04_generation/q05/s01}

\recap{Мазер поддерживает собственное поле за счёт отобранных возбуждённых частиц. Пассивный стандарт использует квантовый переход как частотный дискриминатор внешнего генератора. Порог, ширина, стабильность и систематический сдвиг — разные характеристики.}
\board{Сортировка молекулярного пучка; поле–поляризация; схема петли атомной стабилизации и две зоны Рамзе.}
\oral{Почему точное значение частоты идеального перехода не делает любой атомный прибор безошибочными часами?}
\bait{Можно оставить резонаторное затягивание мазера: частота определяется не только молекулой, но и фазовым откликом резонатора.}
\sources{NIST~\cite{nisttime}, определение секунды и объяснение атомного стандарта; квантовая модель усиления сопоставлена со скоростной моделью вопроса~\Qref{1.10.1}. Присутствующие на Drive видеозаписи стандартов времени целиком не просмотрены; схема мазера изложена как стандартная модель.}
